\SixthArticle{18}
Wir führen von nun an folgende Bezeichnung ein:
\[
[38]\qquad \Pi(k,z)=\frac{1\mathbin{.}2\mathbin{.}3\mathinner{.\,.\,.\,.\,.}k}{(z+1)(z+2)(z+3)\mathinner{.\,.\,.\,.}(z+k)}k^z
\]
wobei unter $k$ seiner Natur nach eine positive ganze Zahl zu verstehen ist; mit dieser Einschränkung stellt $\Pi(k,z)$ eine vollständig bestimmte Funktion der beiden Größen $k,z$ dar. Damit lässt sich leicht erkennen, dass der am Ende des vorigen Artikels aufgestellte Satz folgendermaßen dargestellt werden kann:
\[
[39]\qquad F(\alpha,\beta,\gamma,1)=\frac{\Pi(k,\gamma-1)\mathbin{.}\Pi(k,\gamma-\alpha-\beta-1)}{\Pi(k,\gamma-\alpha-1)\mathbin{.}\Pi(k,\gamma-\beta-1)}\mathbin{.}F(\alpha,\beta,\gamma+k,1)
\]
\SixthArticle{19}
Es wird sich lohnen, die Beschaffenheit der Funktion $\Pi(k,z)$ genauer zu betrachten. Ist $z$ eine negative ganze Zahl, so nimmt die Funktion offenbar einen unendlich großen Wert an, sobald $k$ ein hinreichend großer Wert gegeben wird. Für nichtnegative ganzzahlige Werte von $z$ haben wir dagegen
\[
\begin{aligned}
\Pi(k,0)&=1\\[5pt]
\Pi(k,1)&=\frac{1}{1+\frac1k}\\[5pt]
\Pi(k,2)&=\frac{1\mathbin{.}2}{(1+\frac1k)(1+\frac2k)}\\[5pt]
\Pi(k,3)&=\frac{1\mathbin{.}2\mathbin{.}3}{(1+\frac1k)(1+\frac2k)(1+\frac3k)}
\end{aligned}
\]
usw. oder allgemein
\[
[40]\qquad \Pi(k,z)=\frac{1\mathbin{.}2\mathbin{.}3\mathinner{.\,.\,.\,.\,.\,.}z}{(1+\frac1k)(1+\frac2k)(1+\frac3k)\mathinner{.\,.\,.\,.}(1+\frac zk)}
\]
\SourceLeaf{154}
Allgemein haben wir aber für \textit{jeden beliebigen} Wert von $z$
\[
[41]\qquad \Pi(k,z+1)=\Pi(k,z)\mathbin{.}\frac{1+z}{1+\frac{1+z}{k}}
\]
\[
[42]\qquad \Pi(k+1,z)=\Pi(k,z)\mathbin{.}\left\{\frac{(1+\frac1k)^{z+1}}{1+\frac{1+z}{k}}\right\}
\]
und somit, da $\Pi(1,z)=\frac1{z+1}$ ist,
\[
[43]\qquad \Pi(k,z)=\frac1{z+1}\mathbin{.}\frac{2^{z+1}}{1^z\mathbin{.}(2+z)}\mathbin{.}\frac{3^{z+1}}{2^z(3+z)}\mathbin{.}\frac{4^{z+1}}{3^z(4+z)}\mathinner{.\,.\,.\,.}\frac{k^{z+1}}{(k-1)^z(k+z)}
\]
\SixthArticle{20}
Besondere Aufmerksamkeit verdient jedoch der \textit{Grenzwert}, dem sich für einen gegebenen Wert von $z$ die Funktion $\Pi(k,z)$ stetig nähert, während $k$ ins Unendliche wächst. Es sei zunächst $h$ ein endlicher Wert von $k$, der größer als $z$ ist. Dann ist klar, dass, wenn man $k$ von $h$ zu $h+1$ übergehen lässt, der Logarithmus von $\Pi(k,z)$ einen Zuwachs erhält, der durch folgende konvergente Reihe ausgedrückt wird:
\[
\frac{z(1+z)}{2(h+1)^2}+\frac{z(1-zz)}{3(h+1)^3}+\frac{z(1+z^3)}{4(h+1)^4}+\frac{z(1-z^4)}{5(h+1)^5}+\quad\text{usw.}
\]
Geht also $k$ vom Wert $h$ zu $h+n$ über, so erhält der Logarithmus von $\Pi(k,z)$ den Zuwachs
\[
\begin{aligned}
&\tfrac12 z(1+z)\left(\frac1{(h+1)^2}+\frac1{(h+2)^2}+\frac1{(h+3)^2}+\text{usw.}+\frac1{(h+n)^2}\right)\\[5pt]
&+\tfrac13 z(1-zz)\left(\frac1{(h+1)^3}+\frac1{(h+2)^3}+\frac1{(h+3)^3}+\text{usw.}+\frac1{(h+n)^3}\right)\\[5pt]
&+\tfrac14 z(1+z^3)\left(\frac1{(h+1)^4}+\frac1{(h+2)^4}+\frac1{(h+3)^4}+\text{usw.}+\frac1{(h+n)^4}\right)\\[5pt]
&+\text{usw.}
\end{aligned}
\]
von dem sich leicht beweisen lässt, dass er stets endlich bleibt, selbst wenn $n$ ins Unendliche wächst. Sofern also nicht schon in $\Pi(h,z)$ ein unendlicher Faktor vorhanden war, das heißt, sofern $z$ keine negative ganze Zahl ist, wird der Grenzwert von $\Pi(k,z)$ für $k=\infty$ gewiss eine endliche Größe sein. Offenbar hängt somit $\Pi(\infty,z)$ nur von $z$ ab oder stellt eine vollständig bestimmte Funktion von $z$ dar, die von nun an einfach mit $\Pi z$ bezeichnet wird. Wir definieren daher die Funktion $\Pi z$ durch den Wert des Produkts
\[
\frac{1\mathbin{.}2\mathbin{.}3\mathinner{.\,.\,.\,.\,.}k\mathbin{.}k^z}{(z+1)(z+2)(z+3)\mathinner{.\,.\,.\,.}(z+k)}
\]
\SourceLeaf{155}
für $k=\infty$ oder, wenn man will, durch den Grenzwert des unendlichen Produkts
\[
\frac{1}{z+1}.\frac{2^{z+1}}{1^z(2+z)}.
\frac{3^{z+1}}{2^z(3+z)}.\frac{4^{z+1}}{3^z(4+z)}\quad\mathrm{usw.}
\]

\Needspace{24\baselineskip}
\SixthArticle{21}
Aus Gleichung 41 folgt unmittelbar die Grundgleichung
\[
\begin{array}{@{}ll@{}}
[44]&\displaystyle\Pi(z+1)=(z+1)\Pi z
\end{array}
\]
und daraus allgemein, wenn $n$ eine beliebige positive ganze Zahl bezeichnet,
\[
\begin{array}{@{}ll@{}}
[45]&\displaystyle\Pi(z+n)=(z+1)(z+2)(z+3)\,.\,.\,.\,.(z+n)\Pi z
\end{array}
\]
Für einen negativen ganzzahligen Wert von $z$ ist der Wert der Funktion $\Pi z$ unendlich groß; für nichtnegative ganzzahlige Werte haben wir
\[
\begin{array}{rcl}
\Pi0&=&1\\
\Pi1&=&1\\
\Pi2&=&2\\
\Pi3&=&6\\
\Pi4&=&24\quad\mathrm{usw.}
\end{array}
\]
und allgemein
\[
\begin{array}{@{}ll@{}}
[46]&\displaystyle\Pi z=1.2.3\,.\,.\,.\,.z
\end{array}
\]
Es ist aber verfehlt, diese Eigenschaft unserer Funktion als ihre Definition auszugeben: Sie ist ihrer Natur nach auf ganzzahlige Werte beschränkt und kommt außer unserer Funktion noch unendlich vielen anderen zu (zum Beispiel $\cos2\pi z.\Pi z$, $\cos\pi z^{2n}\Pi z$ usw., wobei $\pi$ den halben Umfang eines Kreises bezeichnet, dessen Radius $=1$ ist).

